\chapter{Every primitive square survives}\label{ch:even}
\section{The local numerator introduced by content}
The transform in \eqref{eq:mobius} is a Dirichlet convolution of multiplicative functions, so $R$ is multiplicative. At a prime $p$, write
\[
a=\tau(p),\qquad b=p^{11},\qquad c=\chifive(p)p^{10}.
\]
The local generating factor is
\begin{equation}\label{eq:rlocal}
\boxed{\sum_{r\ge0}R(p^r)X^r
=\frac{1-cX^2}{1-aX+bX^2}.}
\end{equation}
It follows either by multiplying the Hecke factor by the square-supported Möbius factor, or by reading the two possible square divisors with nonzero Möbius value at a prime power. In particular,
\[
R(1)=1,\qquad R(p)=a,\qquad
R(p^r)=T_r-cT_{r-2}\quad(r\ge2).
\]
The subtraction begins at $r=2$. No undefined $T_{-1}$ or $T_{-2}$ is needed.

The even primitive blocks are not the same as the even tau blocks. For example,
\[
R(p^2)=a^2-p^{10}(p+\chifive(p)).
\]
An argument that $\tau(p^{2m})\ne0$ does not alone imply the corresponding difference is nonzero. We now prove the primitive statement directly.

\begin{theorem}\label{thm:primitiveeven}
For every prime $p$ and every integer $m\ge1$,
\[
R(p^{2m})\ne0.
\]
Consequently $R(s^2)\ne0$ for every positive integer $s$.
\end{theorem}

The proof is stronger than the actual-tau specialization in a useful way. For $p\ne5$, it works for every integer recurrence parameter $a$ with $a\equiv p+\chifive(p)\pmod5$. At $p=5$, it works for every integer $a$. It therefore includes a hypothetical parameter $a=0$ whenever that parameter is compatible with the congruence. The square theorem has not secretly assumed prime nonvanishing.

\section{A residue eliminates three prime classes}
Suppose $p\ne5$ and put $\chi=\chifive(p)$. Modulo $5$, $\chi\equiv p^2$, so the two recurrence roots may be taken as $p$ and $\chi$:
\[
p+\chi\equiv a,\qquad p\chi\equiv p^{11}\pmod5.
\]
The complete homogeneous recurrence coefficient is
\[
T_r\equiv C_r=\sum_{i=0}^r p^{r-i}\chi^i\pmod5.
\]
Also $\chi p^{10}\equiv1\pmod5$. Since $\chi^2=1$, for $r\ge2$ the sum telescopes:
\begin{equation}\label{eq:rmod5}
R(p^r)\equiv C_r-C_{r-2}
=p^{r-1}(p+\chi)\pmod5.
\end{equation}
The same expression holds at $r=1$. It is nonzero for $p\equiv1,2,3\pmod5$. It proves nonvanishing for all positive exponents in those prime classes. Only $p=5$ and $p\equiv4\pmod5$ remain.

\section{The recurrence polynomial and untied valuations}
For $p\equiv4\pmod5$, $\chi=1$ and $a\equiv0\pmod5$. The exact recurrence expansion is
\begin{equation}\label{eq:trexpansion}
T_r=\sum_{j=0}^{\lfloor r/2\rfloor}
(-1)^j\binom{r-j}{j}b^ja^{r-2j}.
\end{equation}
It follows by induction: multiplying the previous row by $a$ and subtracting $b$ times the row before it combines neighboring binomial coefficients. Equivalently, expand the rational generating factor formally. The extreme coefficients are $1$, so a strictly least valuation at either end cannot be raised by a coefficient divisible by $p$.

First suppose $a\ne0$, and write $t=v_p(a)$. In $T_{2m}$, the term indexed by $j$ has valuation at least
\[
2mt+j(11-2t).
\]
If $t\le4$, the $j=0$ term $a^{2m}$ has strictly least valuation $2mt$. Likewise $T_{2m-2}$ has valuation $(2m-2)t$. The correction $p^{10}T_{2m-2}$ has valuation
\[
2mt+(10-2t)>2mt.
\]
Thus
\begin{equation}\label{eq:eventlow}
v_p(R(p^{2m}))=2mt,\qquad t\le4.
\end{equation}

If $t\ge6$, the constant-in-$a$ term $(-b)^m$ is uniquely least-valued in $T_{2m}$, with valuation $11m$. The analogous term gives $v_p(T_{2m-2})=11(m-1)$. The correction has valuation $11m-1$, now smaller than the first block. Hence
\begin{equation}\label{eq:eventhigh}
v_p(R(p^{2m}))=11m-1,\qquad t\ge6.
\end{equation}
Actual nonzero tau parameters at $p\ge5$ cannot have $t\ge6$ by Deligne's bound. The wider argument is nonetheless a useful exact control for the recurrence class.

If $a=0$, do not assign it a finite $t$. Direct substitution gives
\begin{equation}\label{eq:evenzeroa}
R(p^{2m})=(-1)^mp^{11m-1}(p+1),
\end{equation}
which is nonzero. This branch is indispensable: a theorem purporting to survive a hypothetical prime zero must actually evaluate that case.

\section{The tied valuation \texorpdfstring{$t=5$}{t=5}}
The preceding endpoint arguments leave $t=5$. Write $a=p^5u$ with $p\nmid u$, and set
\[
H_r(u)=p^{-5r}T_r.
\]
Equation \eqref{eq:trexpansion} shows that $H_r$ is an integral polynomial. Its recurrence is
\[
H_0=1,\qquad H_1=u,\qquad H_r=uH_{r-1}-pH_{r-2}.
\]
The primitive block becomes
\begin{equation}\label{eq:tied}
R(p^{2m})=p^{10m}\bigl(H_{2m}(u)-H_{2m-2}(u)\bigr).
\end{equation}
The expression in parentheses is a monic polynomial $F_m(x)$ in $x=u^2$. Its constant term and reduction modulo $p$ are
\begin{equation}\label{eq:fm}
F_m(0)=(-1)^mp^{m-1}(p+1),\qquad
F_m(x)\equiv x^{m-1}(x-1)\pmod p.
\end{equation}
The leading term comes from $u^{2m}$. The constant term comes from $H_{2m}(0)=(-p)^m$ and $H_{2m-2}(0)=(-p)^{m-1}$.

Suppose $R(p^{2m})=0$. Then $x=u^2$ is a positive integer root of the monic integer polynomial $F_m$. The integer-root theorem gives $x\mid p^{m-1}(p+1)$. Because $p\nmid u$, one obtains $x\mid p+1$. Reduction in \eqref{eq:fm} gives $x\equiv1\pmod p$. Since $1\le x\le p+1$, the only possibilities are
\[
u^2=1\quad\text{or}\quad u^2=p+1.
\]
The first is impossible: $a$ is divisible by $5$ and $p$ is not, so $u$ is divisible by $5$. The second gives
\[
p=(|u|-1)(|u|+1).
\]
For a prime this forces $|u|=2$ and $p=3$, outside the family $p\equiv4\pmod5$. Both possibilities are excluded. This closes the tied case for every exponent $m$, rather than for a tested list of exponents.

For $m=1$, the polynomial is simply $F_1(x)=x-(p+1)$. For larger $m$, its degree grows, but the same constant-term and residue argument controls any integral square root. That is why no bound on $m$ is needed.

\section{The prime 5 and multiplication over squares}
At $p=5$, $\chifive(5)=0$, so $R(5^r)=T_r$. If $a=0$, its even blocks are $(-5^{11})^m$. If $a\ne0$, then $2v_5(a)$ is even and cannot equal $11$. In the expansion \eqref{eq:trexpansion}, either the first term or the constant-in-$a$ term has a unique least valuation at an even exponent. Thus no even block vanishes.

All prime cases have now been covered. If $s=\prod p^{m_p}$, multiplicativity gives
\[
R(s^2)=\prod_pR(p^{2m_p}).
\]
Every factor is nonzero, and $R(1)=1$. This proves Theorem~\ref{thm:primitiveeven}.

The theorem gives survival, not a common sign. For example,
\[
R(4)=-448,\qquad R(9)=-54594,\qquad R(25)=-25499225,
\]
while other square values can have different signs. Nor does it furnish a fixed auxiliary-prime unit: a nonzero integer can be divisible by arbitrarily high preassigned moduli at suitable varying primes. Chapter~\ref{ch:congruences} makes that last observation into an actual-prime theorem.

\figplaceholder{square_residuals}{Exact primitive residuals at square indices $s^2$ for $1\le s\le22$, displayed on a signed logarithmic scale. Their survival is a consequence of the universal theorem, whose proof also allows a compatible zero prime parameter. The graph illustrates the theorem rather than proves it.}

\chapter{Odd primitive blocks retain exactly the prime obstruction}\label{ch:odd}
\section{Polynomial division before evaluation}
The odd Hecke polynomial $T_{2m+1}$ is divisible by $a$. For $m\ge1$, so is
\[
T_{2m+1}-cT_{2m-1}.
\]
Define its quotient as an integer polynomial in $a^2$ before substituting the actual value $a=\tau(p)$. This matters at a hypothetical zero: the quotient exists algebraically even when numerical division by $a$ would be undefined.

\begin{theorem}\label{thm:oddmultiplier}
For every odd prime $p\ne5$ and every $m\ge0$, an odd integer $U_{p,m}$ satisfies
\[
R(p^{2m+1})=\tau(p)U_{p,m}.
\]
The multiplier is defined by the recurrence polynomial, including when $\tau(p)=0$. If $\tau(p)\ne0$, then
\[
v_2(R(p^{2m+1}))=v_2(\tau(p)).
\]
\end{theorem}
\begin{proof}
For $m=0$, take $U_{p,0}=1$. For $m\ge1$, use the quotient
\[
U(a^2)=\frac{T_{2m+1}-cT_{2m-1}}a
\]
as polynomial division. The constant term in $a^2$ comes from the linear-in-$a$ terms in the two blocks. By \eqref{eq:trexpansion}, it is
\begin{equation}\label{eq:oddconstant}
(-1)^mb^{m-1}\bigl((m+1)b+mc\bigr).
\end{equation}
For an odd prime $p\ne5$, both $b=p^{11}$ and $c=\chifive(p)p^{10}$ are odd. The bracket in \eqref{eq:oddconstant} is congruent to $(m+1)+m=2m+1$ modulo $2$. Thus the constant term is odd.

The actual $a=\tau(p)$ is even by \eqref{eq:parity}. Every nonconstant term in the polynomial in $a^2$ is even after evaluation. Hence the evaluated quotient is odd. It cannot be zero. Multiplication by $a$ gives the theorem, and multiplying a nonzero integer by an odd integer preserves its $2$-adic valuation.
\end{proof}

This proof uses parity rather than a delicate size estimate for the quotient. It also exhibits the zero parameter honestly: if $a=0$, every odd primitive block is zero and every even primitive block remains nonzero by Chapter~\ref{ch:even}.

\section{The primes 2 and 5}
At $p=2$, the residue formula \eqref{eq:rmod5} gives
\[
R(2^r)\equiv2^{r-1}(2-1)\not\equiv0\pmod5
\]
for every $r\ge1$. Thus no local block at $2$ vanishes.

At $p=5$, the character is zero and $R(5^r)=\tau(5^r)$. The actual parameter is $\tau(5)=4830$, with $v_5(a)=1$. In \eqref{eq:trexpansion}, the $j=0$ term has valuation $r$ and every other term has valuation at least $r+9j$. Hence
\[
v_5(R(5^r))=r
\]
for every positive exponent, including odd ones.

\section{The exact zero-equivalence theorem}
\begin{theorem}\label{thm:zeroequivalence}
For every positive integer $n$,
\[
R(n)=0\quad\Longleftrightarrow\quad\tau(n)=0.
\]
More precisely, either side is zero exactly when some prime occurring to an odd exponent in $n$ has $\tau(p)=0$.
\end{theorem}
\begin{proof}
The primitive transform is multiplicative. Its even prime-power blocks are nonzero by Theorem~\ref{thm:primitiveeven}. Its odd blocks at odd primes other than $5$ are the prime parameter times an odd, hence nonzero, integer by Theorem~\ref{thm:oddmultiplier}. The exceptional local blocks at $2$ and $5$ are nonzero by the preceding calculations. Thus its product over the prime factorization of $n$ has precisely the claimed zero criterion. Theorem~\ref{thm:primereduction} gives the same criterion for $\tau(n)$.
\end{proof}

Primitive extraction introduces no new zero. It also removes no zero of the original coefficient sequence. This is stronger and more informative than agreement at a finite list of shells. It describes exactly how the two arithmetic functions encode the same conjecture.

\section{The theorem's gain and its limit}
The primitive square theorem and the odd multiplier theorem classify the new local factors. They show that the extraction does not add an accidental vanishing caused by the numerator $1-\chifive(p)p^{10}X^2$. They also give a precise dyadic valuation law at every odd nonzero block.

The remaining obstruction is concentrated in the first prime parameter. For $p\equiv4\pmod5$, the hypotheses used in the even theorem permit $a=0$. The resulting blocks are
\[
\begin{aligned}
R(p^{2m})&=(-1)^mp^{11m-1}(p+1),&&m\ge1,\\
R(p^{2m+1})&=0,&&m\ge0.
\end{aligned}
\]
The constant block is $R(1)=1$. This is consistent with the congruence, parity and recurrence inputs. Therefore those inputs and the new square theorem cannot jointly establish $a\ne0$. They have isolated its role.

At a prime shell there is no imprimitive part, so the universal target is still
\[
\sum_{\substack{v\in\mathcal V_p\\g(v)=1}}W(v)
=R(p)=\tau(p)\ne0.
\]
The content geometry has clarified the radial arithmetic. The prime total still needs a global or uniform arithmetic mechanism excluding cancellation to zero.
